\chapter{人机交互}

\section{领域导言与演进背景}

\begin{preamblebox}[知识领域概述与变革]
计算系统不仅赋予用户解决现实复杂问题的能力，更促进了人与人之间的社会连接，全方位赋能人类多元化的实践与探索。因此，这些系统在与用户交互及解决问题时，必须恪守尊重个体尊严、维护社会正义、捍卫人类崇高价值与激发创造力的基本准则。人机交互 (HCI)正是一门从跨学科视角切入的核心领域，深度融汇了认知心理学、人因工程学、商业战略与交互设计原则。

每一个用户都是独一无二的。从人机交互的视角出发，任何与人类交互的系统设计，都必须前瞻性地预见并充分尊重这种多样性。这不仅涵盖无障碍辅助接入，更延伸至文化习俗与社会规范、神经多样性、多感官交互模态以及系统在用户心理与生理上激发的全方位反应。一个高效卓越的计算系统，应当在公平对待用户、严守隐私边界、构筑安全屏障并恪守伦理准则的同时，赢得用户的深层信任。

实现上述目标，必须确立以“设计为中心”的工程范式。该范式始于明确的设计意图，并深刻认识到设计是一个动态的迭代过程——需要针对系统的可用性及其对最终用户的深远影响进行持续且反复的评测与验证。此外，现代技术不仅通过其计算产出激发用户反馈，其感知环境与交互通信的多样化模态同样深刻塑造着人机体验。本知识领域致力于系统性提升计算机从业者对这些核心议题的敏感度与工程能力，理应深刻启迪与塑造每一位未来的计算机科学家与工程师。

\textbf{自 CS2013 以来的核心变革：}
\begin{enumerate}
    \item \textbf{知识单元重构与现代实践全面对接}：重新命名并系统重构了各知识单元，精准反映当前产业界面向敏捷交付、用户体验（UX）研究与多端融合的最新工程实践，并前瞻预见空间计算与生成式 AI 等前沿技术的演进趋势。
    \item \textbf{强化多样性本质与以用户为中心的设计}：更加突出用户群体在生理、认知与文化上的内在多样性本质，将“以用户为中心的设计 (UCD)”和“包容性设计”置于计算机专业教育的核心枢纽地位。
    \item \textbf{交互模态视角的现代化演进}：多模态交互（如文本、图形、语音、触觉与姿态）依然占据核心地位，但大纲降低了对孤立特定输入设备的机械强调，转而倡导更加与时俱进、富有同理心与情境感知的全通道交互工程方法。
    \item \textbf{深度关注系统社会影响与伦理可解释性}：课程体系高度强调系统评估计算软硬件对用户的长远影响，将可持续性、算法公平性、防范差别待遇、系统透明度、可信度与可解释性贯穿于交互全生命周期。
    \item \textbf{确立跨学科交叉领域的枢纽定位}：鉴于与计算机学科各分支的广泛交织，HCI 本身即是一个核心交叉领域，与人工智能（AI）、社会、伦理与专业素养（SEP）、软件开发基础（SDF）以及软件工程（SE）等知识领域形成了深度的协同互补。
\end{enumerate}
\end{preamblebox}

\section{核心学时分配建议}

本知识领域在“向日葵模型”下的学时推荐如下表所示：

\begin{table}[H]
\centering
\setlength{\tabcolsep}{3.5pt}\fontsize{8pt}{9.6pt}\selectfont\begin{tabularx}{\columnwidth}{Y c c}
\toprule
\textbf{知识单元} & \textbf{计算机核心} & \textbf{领域核心} \\
\midrule
理解用户 & 2 & 5 \\
设计中的问责与责任 & 2 & 2 \\
无障碍与包容性设计 & 2 & 2 \\
交互设计评估 & 1 & 2 \\
交互系统设计 & 1 & 5 \\
\midrule
\textbf{合计} & \textbf{8 学时} & \textbf{16 学时} \\
\bottomrule
\end{tabularx}
\end{table}

\noindent\textbf{说明}：知识单元 6（人机交互与社会伦理，HCI and SEP）属于贯穿整个计算机科学培养体系的跨领域交叉单元，其核心学时已与《社会、伦理与专业素养 (SEP)》深度有机融合，并在本领域各单元教学实践与期末综合项目中持续渗透。

\section{知识单元与详细大纲}

\subsection{理解用户：个体目标与人际交互}
\textbf{【学时建议】}：2 计算机核心学时，5 领域核心学时。

本单元奠定人本设计的核心认知基石。学生必须打破以自我为中心的工程师思维惯性，深刻理解真实用户的生理机能、认知模型与多元社会背景，掌握科学的用户需求挖掘与画像建模方法。

\subsubsection{计算机核心主题}
\begin{itemize}
    \item \textbf{以用户为中心的设计与评估方法}：
    \begin{itemize}
        \item “你不是用户”的核心哲学内涵及其对软件工程思维的重塑。
        \item 用户需求挖掘技术：探索性调研、定性观察与用户痛点定位。
        \item 形成性研究：在概念设计与早期原型阶段探寻需求规律。
        \item 核心调研手段：结构化/半结构化深度访谈、问卷调查法、基线可用性实测。
    \end{itemize}
\end{itemize}

\subsubsection{领域核心主题}
\begin{itemize}
    \item \textbf{用户中心设计方法论}：
    \begin{itemize}
        \item 用户画像构建与画像图谱：典型用户与极端边界用户的特征抽象与共情载体。
        \item 用户故事与叙事构建：敏捷用户故事模板、场景化故事板与叙事采集方法。
        \item 同理心地图：说（Says）、想（Thinks）、做（Does）、感受四象限建模。
        \item 进阶需求评估技术：民族志观察与情境调查。
        \item 用户体验旅程图：触点分析、情绪曲线映射与痛点识别（参见“交互设计评估”单元）。
    \end{itemize}
    \item \textbf{用户的生理机能与认知特征}：
    \begin{itemize}
        \item 指导交互设计的生理学基础：人机工程学、色彩感知生理机制（视网膜杆体/锥体细胞特性与色盲模拟）。
        \item 指导交互设计的认知心理学模型：注意力分配机制（选择性/持续性注意）、知觉感知与图形识别机制、运动控制模型（菲茨定律 Fitts' Law 初步）、人类记忆多存储模型（感觉记忆、工作记忆容量限制 $7\pm2$ 效应、长时记忆提取）。
        \item 社会与行为心理学课题：认知偏差（Confirmation Bias、锚定效应）、变化盲视与界面诱导心理。
    \end{itemize}
    \item \textbf{面向多元用户群体的设计}：
    \begin{itemize}
        \item 用户属性多元性：种族、生理能力、年龄阶层、性别、文化习俗、数字素养经验与受教育程度对交互体验与用户期望的决定性影响。
        \item 软件国际化 (Internationalization, i18n) 与跨文化设计：跨文化维度理论在界面布局、隐喻及信息架构中的映射。
        \item 有效设计评估面临的严峻挑战：被试抽样偏差、实验结论的普适泛化性局限；残障人士与功能受限者的真实交互体验感知。
        \item 通用设计哲学（参见“无障碍与包容性设计”单元）。
    \end{itemize}
    \item \textbf{协同交互与群体通信}：
    \begin{itemize}
        \item 多用户与组织情境下的用户认知特征：角色分工、权力结构与群组动力学。
        \item 同步群体通信系统：实时聊天、多方音视频会议协同、多人在线协作虚拟空间。
        \item 异步群体通信机制：电子邮件线程、主题论坛、协同文档版本流转、社交网络。
        \item 社交媒体、社交计算 与社交网络分析 (SNA) 初步。
        \item 在线协同、社会化协调机制与虚拟数字社区治理。
        \item 虚拟化身、数字化身交互与虚拟三维空间协同。
    \end{itemize}
\end{itemize}

\subsubsection{拓展选修与前沿主题}
\begin{itemize}
    \item 群体智能与协同众包计算系统架构；元宇宙沉浸式数字空间中具身虚拟化身的非语言交流与微表情映射。
\end{itemize}

\subsubsection{示范性学习成果与能力目标}
\begin{enumerate}
    \item \textbf{[应用]} 在一个真实工程软件开发项目中，规范组织并端到端实施“以用户为中心的设计（UCD）”全流程。
    \item \textbf{[解释]} 深入对比与辨析最终用户真实诉求与系统软件工程师/架构师先验技术偏见之间的本质分歧。
    \item \textbf{[解释]} 精准识别给定系统的代表性目标用户画像，并系统分析该系统可能对周边间接利益相关者产生的连锁影响。
    \item \textbf{[解释]} 阐明在现代交互设计全流程中，“同理心构建”与“科学评估”作为核心驱动要素的方法论意义。
    \item \textbf{[应用]} 独立组织针对真实目标用户的调研工作，并撰写符合工业规范的用户特征与多维度需求分析报告。
    \item \textbf{[开发]} 基于系统化的需求评估结果，归纳提炼出满足敏捷开发标准的高质量用户故事。
    \item \textbf{[开发]} 针对特定初始受众设计的产品方案进行可用性重构，使其精准适配文化背景或认知能力存在显著差异的新用户群体。
    \item \textbf{[评估]} 针对给定的实际设计难题，横向权衡并对比不同需求挖掘方法的适用场景、成本开销与信息捕获深度。
    \item \textbf{[评估]} 批判性反思并评估当前系统方案中引入低技术或无技术非数字组件所能带来的系统增益与体验改善。
    \item \textbf{[解释]} 深刻剖析多用户协作系统与社交计算情境对交互架构、并发响应与冲突消解带来的深层技术影响。
\end{enumerate}

\subsection{设计中的问责与责任：可持续性、安全、隐私、信任与伦理}
\textbf{【学时建议】}：2 计算机核心学时，2 领域核心学时。

软件交互并非价值中立。交互设计师与软件工程师做出的每一个界面与逻辑决策，均深刻影响着资源消耗、个人隐私、社会公平及人类身心健康。本单元引导学生建立负责任的工程伦理与问责意识。

\subsubsection{计算机核心主题}
\begin{itemize}
    \item \textbf{交互设计的系统性深远影响}：
    \begin{itemize}
        \item 绿色低碳与可持续性：算力与带宽感知的轻量化界面设计、减少数字废弃物与设备过度消耗。
        \item 社会包容性 与系统安全性 (Safety \& Security)：防范界面诱导导致的误操作与财产/人身安全风险。
        \item 数据隐私 与防范技术危害：知情同意权界定、敏感数据非必要不采集原则。
        \item 差别待遇与歧视性影响：算法推荐与界面交互对特定群体的隐性排斥。
    \end{itemize}
    \item \textbf{交互设计中的职业伦理}：
    \begin{itemize}
        \item 设计方法与工程解决方案中的伦理底线；人工智能技术介入交互系统（如生成式 AI、自主决策代理）时的边界规范。
        \item 统筹考量利益相关者多维影响与人因工程的责任担当。
        \item 坚持以人为本、以解决用户真实痛点为导向的设计使命，坚决抵制损害用户权益的病态 KPI 驱动设计。
    \end{itemize}
    \item \textbf{工程设计中的法规与合规约束}：
    \begin{itemize}
        \item 软件所有权归属与问责边界。
        \item 法律法规合规框架：通用数据保护条例（GDPR）、个人信息保护法、未成年人网络保护条例等合规基准。
        \item 超越用户即时表层欲望的长效考量：通过对工程问题本质的持续反思与迭代重构，防范成瘾机制与负面社会外部性。
    \end{itemize}
\end{itemize}

\subsubsection{领域核心主题}
\begin{itemize}
    \item \textbf{价值敏感设计}：
    \begin{itemize}
        \item 直接利益相关者 与间接利益相关者 的精准辨识与诉求映射。
        \item 系统梳理多元利益相关者的价值体系，建立兼顾尊严、自由、公正与信任的工程价值权衡模型。
    \end{itemize}
    \item \textbf{界面诱导与说服性设计}：
    \begin{itemize}
        \item 科学量化与评估交互界面中的说服性元素。
        \item 正向健康行为引导（如促进运动、助益专注）作为合法合规的设计目标。
        \item 批判性识别与抵制“暗黑模式”：诱导式收费、隐藏取消订阅通道、虚假紧迫感及情感绑架等恶意交互手段。
    \end{itemize}
\end{itemize}

\subsubsection{示范性学习成果与能力目标}
\begin{enumerate}
    \item \textbf{[评估]} 敏锐识别并批判性审视系统设计方案对社会及目标社群潜在的多维冲击，系统应对可持续性、包容性、安全性、隐私保护、潜在危害及差别对待等重大现实关切。
    \item \textbf{[解释]} 准确识别与剖析给定工程设计中潜在的人因工程学要素。
    \item \textbf{[解释]} 系统梳理并深刻洞察产品生态中的直接利益相关者与间接利益相关者群体。
    \item \textbf{[开发]} 拓展构建覆盖产品完整生命周期的情境剧本，前瞻预判超出初始设计意图之外的长周期社会与群体演变。
    \item \textbf{[评估]} 批判性审视系统设计中直接或间接影响利益相关者与更广泛社会的决定性因素（如算法透明度、系统可持续生存力、用户信任机制、AI 自主性边界等）。
    \item \textbf{[评估]} 严格审视产品界面与交互逻辑中的说服性元素及其设计初衷，客观评估其是否与用户的长远真实利益保持一致。
    \item \textbf{[评估]} 对照系统初始设计意图，客观评估系统在真实环境下的最终实施成效与衍生产出。
    \item \textbf{[解释]} 深刻领会关键架构与交互设计决策在社会、组织与个体层面引发的连锁反应与长远影响。
\end{enumerate}

\subsection{无障碍与包容性设计}
\textbf{【学时建议】}：2 计算机核心学时，2 领域核心学时。

信息无障碍是现代数字文明与人权保障的关键基石。计算机专业毕业生必须树立“无障碍设计使所有人受益”的崇高理念，精通国际主流无障碍规范标准，熟练掌握支持各类机能受限群体的辅助开发技术。

\subsubsection{计算机核心主题}
\begin{itemize}
    \item \textbf{社会责任与法律基准}：
    \begin{itemize}
        \item 保障残障群体数字平权的法律法规支撑与强制性义务（如中国《无障碍环境建设法》、美国 ADA 法案、欧盟欧洲无障碍法案 EAA）。
        \item 核心人本理念：“无障碍设计使所有人受益”——解决极端困难场景的设计往往能极大提升常规场景的用户体验。
    \end{itemize}
    \item \textbf{无障碍工程标准与技术规范}：
    \begin{itemize}
        \item 万维网联盟（W3C）Web 内容无障碍指南（WCAG 2.1/2.2 AA 级）核心技术原则：可感知性、可操作性、可理解性与健壮性。
    \end{itemize}
    \item \textbf{辅助技术与软硬件赋能产品}：
    \begin{itemize}
        \item 现代操作系统与浏览器内嵌的无障碍功能与辅助技术 (AT)：可访问性对象模型（AOM）、屏幕阅读器、高对比度模式、单手操作模式与物理开关控制。
    \end{itemize}
    \item \textbf{包容性设计框架}：
    \begin{itemize}
        \item 承认并正视个体能力维度的客观差异；通用设计七大核心原则。
    \end{itemize}
\end{itemize}

\subsubsection{领域核心主题}
\begin{itemize}
    \item \textbf{人口统计学与多元障碍形态谱系}：
    \begin{itemize}
        \item 三大典型机能障碍形态：永久性残障（Permanent，如全盲或截瘫）、暂时性机能受损（Temporary，如腕骨骨折或白内障术后）、情境性受限（Situational Disability，如强阳光下屏幕眩光或嘈杂环境中听觉受阻）。
        \item 残障认知的全球化视野与跨文化认知演进；消除技术研发中对特殊群体的隐性歧视与社会偏见。
    \end{itemize}
    \item \textbf{包容性用户体验研究与软件工程实践}：
    \begin{itemize}
        \item 邀请障碍人士参与的全流程包容性 UX 调研与测试方法。
        \item 前端与客户端工程层面的无障碍落地实践：WAI-ARIA 规范标签语义化注入、焦点顺序管理、键盘完全可达性导航、颜色对比度（$\ge 4.5:1$）自动化校验。
    \end{itemize}
    \item \textbf{先进辅助功能与适配技术}：
    \begin{itemize}
        \item 屏幕阅读器严格适配（如 NVDA、JAWS、TalkBack、VoiceOver）；图片多模态替代文本语义化生成；动态字体无级缩放与响应式重排防溢出处理。
    \end{itemize}
    \item \textbf{包容性协同流程与大社会效益}：
    \begin{itemize}
        \item 参与式设计与共同设计流程；以包容性驱动全社会数字鸿沟弥合。
    \end{itemize}
\end{itemize}

\subsubsection{拓展选修与前沿主题}
\begin{itemize}
    \item 认知除蔽与反思性质疑；跨学科残障学理论研究。
    \item 企业实施包容性设计的商业投资回报率量化评估模型。
    \item 用户敏感型包容性设计范式。
    \item 批判性人机交互研究范式：HCI 中的批判性种族理论、女权主义人机交互理论以及批判性残障理论。
\end{itemize}

\subsubsection{示范性学习成果与能力目标}
\begin{enumerate}
    \item \textbf{[解释]} 精准辨识视觉、听觉、肢体与认知等各类机能障碍人士面临的交互瓶颈，明确指出与之相匹配的无障碍与辅助技术解决方案。
    \item \textbf{[解释]} 阐明并应用适宜的包容性设计范式（如通用设计 Universal Design 与基于能力的设计 Ability-Based Design）。
    \item \textbf{[应用]} 熟练掌握并严格遵循主流软件与网络内容无障碍设计指南（如 WCAG 规范族）。
    \item \textbf{[解释]} 准确列举并理解适于无障碍交互软件开发的国内外法律法规与国家强制性标准。
    \item \textbf{[开发]} 灵活将包容性框架运用于视觉界面、语音交互界面（VUI）和触控交互界面的无障碍工程落地，确保全通道无障碍可达。
    \item \textbf{[应用]} 在与多元背景及障碍群体用户进行实地交流访谈与测试时，展现出高度专业、得体且尊重的跨群体沟通素养。
    \item \textbf{[评估]} 综合运用无障碍自动化检测工具与手工辅助技术走查，对现代 Web 页面及移动 App 进行严格的合规审计与无障碍漏洞定位。
    \item \textbf{[评估]} 探究在软硬件工程发展史上，针对残障、种族或性别的隐性偏见如何潜移默化地渗透到技术系统设计中，并从中提炼反思性指导法则。
    \item \textbf{[开发]} 独立构想并实施高阶用户体验研究方案，深度挖掘边缘弱势群体的潜在需求并转化为切实的设计洞察。
\end{enumerate}

\subsection{交互设计评估}
\textbf{【学时建议】}：1 计算机核心学时，2 领域核心学时。

未经科学评估的设计只是未经检验的技术假设。本单元培养学生掌握定性、定量与混合评估范式，能够严谨设计并实施可用性测试实验，借助数理统计工具与质性分析法对交互系统做出可辩护的科学推断。

\subsubsection{计算机核心主题}
\begin{itemize}
    \item \textbf{以用户为中心的评测核心范式}：
    \begin{itemize}
        \item 形成性评估：需求探索、早期纸质原型走查与概念纠偏。
        \item 总结性评估：核心功能验收、交互效率基准测试与可用性达标量化评测。
    \end{itemize}
    \item \textbf{交互系统核心评估维度}：
    \begin{itemize}
        \item 实用性：系统能否真正支撑目标任务的达成。
        \item 操作效率：熟练用户完成既定任务所需耗费的精力、时间与操作击键成本。
        \item 易学性：新手用户初次接触系统达到基线操作水平的学习曲线与门槛。
        \item 用户主观满意度：系统交互引发的用户心理认同与情感愉悦度。
    \end{itemize}
    \item \textbf{伦理审查要求}：
    \begin{itemize}
        \item 在开展涉及人类受试者的可用性测试与数据采集前，必须获取伦理委员会批准（IRB 审查），确保知情同意与隐私脱敏。
    \end{itemize}
\end{itemize}

\subsubsection{领域核心主题}
\begin{itemize}
    \item \textbf{定性、定量与混合研究方法矩阵}：
    \begin{itemize}
        \item 质性分析技术：测试会话录像转录、质性扎根编码与主题提炼分析。
        \item 定量评估方法：任务成功率、任务完成时长、错误率统计与标准心理量表（SUS 系统可用性量表、NASA-TLX 认知负荷量表）。
        \item 混合研究方法：结构化现场观察、出声思考法、深度半结构化访谈与受控对比实验。
        \item 专家走查法：启发式评估（Heuristic Evaluation，重点掌握 Nielsen 十大可用性启发式原则）、认知走查法。
        \item 专业呈现与交付：撰写高质量可用性工程评估白皮书、提炼严谨的用户画像。
        \item 评测方法论瓶颈与挑战：被试抽样偏差、实验生态效度与泛化性限制。
    \end{itemize}
    \item \textbf{可用性实验设计与严谨规划}：
    \begin{itemize}
        \item 明确研究目标，确立科学、可证伪的实验原假设（$H_0$）与备择假设（$H_1$）。
        \item 机构伦理审查委员会（IRB）与伦理合规备案审批全流程。
        \item 开放科学框架下的实验方案预注册机制。
        \item 经典实验设计范式比较：组内设计（Within-Subjects Design，平衡顺序效应与疲劳效应）与组间设计的优劣权衡与控制变量准则。
    \end{itemize}
    \item \textbf{设计产出的综合长效影响评测}：
    \begin{itemize}
        \item 系统对自然生态环境、物质资源消耗、社会结构伦理、网络空间安全与数据隐私构架的深层综合影响评估。
    \end{itemize}
\end{itemize}

\subsubsection{拓展选修与前沿主题}
\begin{itemize}
    \item 定量高级统计分析工具包：R 语言、Python Scipy/Statsmodels、JASP 等在交互数据中的实战；可视分析与探索性数据分析（EDA）。
    \item 统计假设检验进阶：单因素/多因素方差分析（ANOVA）、双样本独立/配对 $t$ 检验、事后多重检验校正（Bonferroni 校正）；参数检验与非参数检验（Wilcoxon 符号秩检验、Mann-Whitney $U$ 检验）的选型条件；效应量（Cohen's $d$、$\eta^2$）计算与统计功效分析。
    \item 评估研究数据合规管理：基于开放科学原则的数据长期安全归档与脱敏共享机制；受试者敏感个人可识别信息（PII）的去标识化与差分隐私注入。
\end{itemize}

\subsubsection{示范性学习成果与能力目标}
\begin{enumerate}
    \item \textbf{[解释]} 深入辨析形成性评估与总结性评估的本质特征、介入节点及其在驱动交互设计持续演进中的关键作用。
    \item \textbf{[应用]} 针对软件系统研发的不同生命周期阶段（从概念草图到发布后监控），精准选取并编排最适宜的形成性或总结性评测方法。
    \item \textbf{[解释]} 深入论证在交互评测中将定性洞察与定量数据相结合的互补优势。
    \item \textbf{[评估]} 针对给定的实际设计方案，全面评估其对社会伦理、环境消耗与用户社群带来的综合宏观影响。
    \item \textbf{[开发]} 独立策划针对特定交互界面的可用性评测实验方案，严谨论证其研究目标设定、实验假设推导及被试控制组设计的合理性。
    \item \textbf{[评估]} 规范组织并执行界面可用性评测，收集多模态实验数据，并在实验范式约束下得出可辩护、高置信度的评估结论。
    \item \textbf{[应用]} 运用专业统计软件对实验收集的行为日志与量表数据执行假设检验，科学解读显著性水平（$p$ 值）与置信区间。
    \item \textbf{[应用]} 完成一份严格符合学术与工业规范的可用性实验设计预注册报告。
\end{enumerate}

\subsection{交互系统设计}
\textbf{【学时建议】}：1 计算机核心学时，5 领域核心学时。

本单元系统阐述从概念草图到工业级软件界面的完整工程构建方法。重点涵盖低保真与高保真原型工具栈、现代前端与客户端软件架构模式、双钻设计流程、格式塔视觉心理学原则以及多通道先进交互技术。

\subsubsection{计算机核心主题}
\begin{itemize}
    \item \textbf{原型设计技术与工具栈}：
    \begin{itemize}
        \item 低保真原型设计：纸质线框图绘制、流程草图、低成本走查与概念试错。
        \item 快速原型制作 与探索废弃型原型。
        \item 原型的交互粒度控制：外观保真度、交互深度与数据动态性的合理平衡。
    \end{itemize}
    \item \textbf{交互系统设计模式}：
    \begin{itemize}
        \item 迭代设计闭环：需求分析 $\to$ 原型构思 $\to$ 实施开发 $\to$ 评估反馈 $\to$ 重构演进。
        \item 通用设计模式与现代交互模式：数据驱动设计与事件驱动交互体系。
    \end{itemize}
    \item \textbf{工程设计约束与权衡}：
    \begin{itemize}
        \item 跨平台异构环境约束（Web、移动端 iOS/Android、桌面端、车载系统、嵌入式终端）。
        \item 设备物理特性限制（屏幕视网膜像素密度、算力能耗配额、触控精度、网络带宽与离线状态）。
    \end{itemize}
\end{itemize}

\subsubsection{领域核心主题}
\begin{itemize}
    \item \textbf{软件架构模式与设计指南}：
    \begin{itemize}
        \item 客户端交互架构：MVC、MVVM、声明式 UI 响应式状态管理（如 React/Flutter 单向数据流）。
        \item 跨平台一致性设计系统：组件库、设计令牌与原子设计理论。
        \item 分布式与多设备协同交互环境下的状态一致性与实时同步挑战。
    \end{itemize}
    \item \textbf{先进设计思维与过程范式}：
    \begin{itemize}
        \item 参与式设计 与共同设计。
        \item 英国设计协会“双钻模型”：探索期发散思维 $\to$ 定义期收敛思维 $\to$ 研发期发散构思 $\to$ 交付期收敛验证。
    \end{itemize}
    \item \textbf{先进交互技术与通道向量}：
    \begin{itemize}
        \item 输入与输出多维通道向量：手势追踪、姿态检测、多点触控、语音指令、力敏反馈。
        \item 现代图形用户界面 (GUI) 精细化组件、物理硬件控制器、触觉反馈技术。
        \item 交互容错机制与优雅降级：防呆设计、异常状态回滚、高容错操作撤销。
    \end{itemize}
    \item \textbf{视觉交互界面设计心理学法则}：
    \begin{itemize}
        \item 色彩理论、对比度阶层与排版网格系统。
        \item 格式塔知觉心理学原理：邻近原则、相似原则、连续原则、闭合原则、主体与背景原则。
    \end{itemize}
\end{itemize}

\subsubsection{拓展选修与前沿主题}
\begin{itemize}
    \item 沉浸式交互空间：虚拟现实（VR）、增强现实（AR）、混合现实（MR）与扩展现实（XR）交互套件开发；空间音频声场定位。
    \item 3D 打印与智能快速成型制造（3D Printing and Digital Fabrication）。
    \item 异步交互协作模型。
    \item 创造力支持工具 (CST)设计原理。
    \item 会话式用户界面（CUI）与生成式智能体对话交互设计。
\end{itemize}

\subsubsection{示范性学习成果与能力目标}
\begin{enumerate}
    \item \textbf{[开发]} 针对特定的应用场景与任务特征，提出并构建量身定制的最优交互模式与系统架构设计方案。
    \item \textbf{[开发]} 严格践行“深度理解用户 $\to$ 增量式原型构建 $\to$ 严密评测反馈 $\to$ 指导下一轮迭代演进”的敏捷交互开发闭环。
    \item \textbf{[解释]} 深入阐述面对动态变化的工程约束（底层硬件变迁、算力限制、用户多元化需求、安全合规收紧）时，系统架构所必须做出的权衡取舍。
    \item \textbf{[评估]} 紧扣核心项目目标，系统评估并抉择主流前端与客户端软件架构模式的优劣。
    \item \textbf{[解释]} 准确定位分布式与多设备协同交互环境中面临的网络并发、状态同步不一致与体验延迟等核心工程挑战。
    \item \textbf{[评估]} 针对特定真实业务情景，深入评估并横向比对不同交互输入技术可能引发的个人数据与生物特征隐私风险。
    \item \textbf{[解释]} 深度运用格式塔等经典视觉认知心理学法则，对特定界面的视觉层级、元素布局与信息架构做出专业的技术原理解释。
    \item \textbf{[评估]} 深入评估空间计算与 VR/AR/MR 虚实融合情境下所暴露的用户空间注视点数据、生物识别及环境空间映射的隐私安全隐患。
\end{enumerate}

\subsection{人机交互与社会、伦理及专业素养的交叉融合}
\textbf{【说明】}：本单元为跨领域核心贯穿单元，与《社会、伦理与专业素养 (SEP)》深度锚定，致力于将人本设计思想与工程师社会责任深度融合，不单独列出静态学时，贯穿于整套课程教学。

\subsubsection{计算机核心主题}
\begin{itemize}
    \item \textbf{通用设计与以用户为中心的设计之伦理意涵}：将普遍可达性、用户尊严维护、反歧视伦理确立为软件工程需求分析与架构设计的第一性原则。
\end{itemize}

\subsubsection{领域核心主题}
\begin{itemize}
    \item \textbf{全流程参与式与包容性设计流程}：保障边缘群体与非技术利益相关者在软件全生命周期中的知情权与决策参与权。
    \item \textbf{交互设计的全方位外部性审计}：系统评测软硬件系统对自然生态环境、物质资源消耗、社会正义、数字安全、数据隐私、职业道德及长远文明构架的系统影响。
\end{itemize}

\subsubsection{示范性学习成果与能力目标}
\begin{enumerate}
    \item \textbf{[应用]} 将以人为本的通用设计规范内嵌于工业软件需求工程与全生命周期开发规范中。
    \item \textbf{[解释]} 深刻辨识计算机工程师在构建大众交互产品时所肩负的法律、伦理与社会正义责任。
    \item \textbf{[评估]} 批判性剖析业界近期发生的缺乏包容性设计或算法偏见的真实工程反面案例，深入论述其引发的恶劣社会后果，并提出针对性的架构整改方案。
    \item \textbf{[评估]} 针对给定的交互软件，系统识别并评估其对生态环境、社会弱势群体与长效文明构建带来的深层外部性影响。
    \item \textbf{[评估]} 全面评估沉浸式虚拟现实（VR/AR/MR）交互情境下的空间隐私与用户数字踪迹保护机制。
\end{enumerate}

\section{专业素养与思维倾向}

在人机交互领域的学习与工程实践中，重点塑造学生的以下专业素养与思维倾向：
\begin{itemize}
    \item \textbf{高度适应性}：在面对快速更迭的软硬件形态、不确定的用户群体以及动态变化的业务需求时，能够敏锐调整设计范式与技术选型。
    \item \textbf{严谨细致的一丝不苟}：在微交互动效、界面边界留白、字体排版层级、无障碍焦点顺序以及异常容错提示等每一个微小工程细节中保持极致追求。
    \item \textbf{深具同理心}：能够超越技术开发者固有的“专家视角”，真正换位思考，体察不同文化背景、不同生理认知条件及边缘群体的真实喜怒哀乐与交互痛点。
    \item \textbf{团队协作导向}：优秀的 HCI 实践者必须将团队的整体成功置于首位，善于在跨职能团队中架起工程师、产品经理、设计师、业务人员与终端用户之间的沟通桥梁。
    \item \textbf{富有创造力}：能够基于历史工程实践、目标受众诉求及 HCI 核心理论，通过非凡的创造力将技术可行性与人本体验完美融合，设计出既优雅又高效的破局方案。
\end{itemize}

\section{数学基础要求}

人机交互领域高度依赖应用数理统计学，以支撑对可用性评估实验数据、A/B 测试结果及海量用户行为日志的科学解析：
\begin{itemize}
    \item \textbf{必须具备的基础统计学核心}：
    \begin{itemize}
        \item 集中趋势度量：均值、中位数、众数在偏态分布与长尾数据下的选用边界。
        \item 离散程度与变异性度量：方差、标准差、四分位距（IQR）与极差。
        \item 频数分布与概率分布规律：正态分布、二项分布与泊松分布；理解长尾分布对用户体验基准测试的影响。
        \item 推断统计基础直觉：参数置信区间、显著性检验（$p$ 值）以及假设检验基本原理。
    \end{itemize}
\end{itemize}

\section{交叉主题与跨学科共享概念}

人机交互作为一门承上启下的枢纽学科，与计算机科学其他关键领域深度交织：
\begin{itemize}
    \item \textbf{与人工智能领域 (AI) 的共享概念}：
    \begin{itemize}
        \item 个人智能助理：多模态自然交互与意图理解。
        \item 机器人学：人机共融协作与人机信任机制。
        \item 增强智能：AI 作为人类智力放大器而非简单替代工具的人机协作范式。
        \item 自主系统中的人工干预与控制权覆盖机制：保障人类核心掌控权的“人在回路”设计。
    \end{itemize}
    \item \textbf{与社会、伦理与专业素养 (SEP) 的共享概念}：
    \begin{itemize}
        \item 数据隐私与行为追踪保护：用户界面中的透明授权与隐私合规。
        \item 数字包容性与社会平权：跨越数字鸿沟与保障弱势群体福祉。
    \end{itemize}
    \item \textbf{与软件工程 (SE) 的共享概念}：
    \begin{itemize}
        \item 产品需求工程与用户故事定义：以人为本的需求引出。
        \item 软件系统验证与可用性确认：形式化测试与真实用户评测的互补。
        \item 与多方利益相关者的协同对接：多角色诉求对齐。
        \item 现代化原型测试、界面协同设计与跨职能沟通工具链。
    \end{itemize}
\end{itemize}

\section{真实能力任务规范}

\begin{taskbox}[任务 1：现代用户界面开发框架横向技术选型与深度权衡对比]
\textbf{【任务情境】}：某企业拟启动大型跨平台数字化业务系统重构，技术委员会需要在主流客户端与前端开发框架（如声明式响应式框架、跨平台自绘引擎与原生平台套件）之间进行工程选型。\\
\textbf{【能力陈述】}：系统阐述特定 UI 框架在应对特定软件工程难题时的优势与局限，清晰证明“界面设计与开发是目标驱动且充满多维权衡”的深刻认知，综合平衡渲染性能、开发维护成本、多端一致性及原生无障碍 API 支持。\\
\textbf{【能力领域】}：软件工程 \quad \textbf{【能力单元】}：需求分析、架构设计、技术文档编制、系统评估与项目管理\\
\textbf{【核心要求】}：计算机核心 \quad \textbf{【达标等级】}：\textbf{评估}
\end{taskbox}

\begin{taskbox}[任务 2：企业核心业务系统可用性研究实验的全流程规划、文档编制与实施]
\textbf{【任务情境】}：某全国性在线医疗问诊挂号平台改版，需针对中老年及视障受试者群体的操作阻塞点，组织严密的可用性实测。\\
\textbf{【能力陈述】}：深刻理解反馈闭环在全生命周期人本设计中的关键驱动作用；根据系统所处的研发阶段与目标受众的认知特征，独立制定兼顾伦理审批、被试招募、任务设计的可用性测试方案，严密组织出声思考走查并产出高质量分析白皮书。\\
\textbf{【能力领域】}：应用工程 \quad \textbf{【能力单元】}：系统设计、技术文档、系统评估、项目管理、用户接受度验证、社会议题调适、系统持续改进\\
\textbf{【核心要求】}：计算机核心 \quad \textbf{【达标等级】}：\textbf{应用与评估}
\end{taskbox}

\begin{taskbox}[任务 3：公共政务系统的可用性、无障碍性与包容性综合审计与改进方案]
\textbf{【任务情境】}：审计某市政“一网通办”数字政务门户，系统排查各类残障人士、高龄老人及数字技能欠缺群体在表单填写与办事指南浏览中的交互障碍。\\
\textbf{【能力陈述】}：熟练运用多元评估透镜（可用性十大原则、W3C WCAG 2.1 AA 级标准与包容性指标），全面诊断面向公众的交互系统在易用性、屏幕阅读器兼容性及排版对比度上的缺陷，制定具有工程可行性的界面重构与代码整改技术方案。\\
\textbf{【能力领域】}：应用工程 \quad \textbf{【能力单元】}：系统评估、技术文档、社会议题调适、系统持续改进\\
\textbf{【核心要求】}：计算机核心 \quad \textbf{【达标等级】}：\textbf{评估}
\end{taskbox}

\begin{taskbox}[任务 4：跨职能敏捷团队协同推进智能交互系统的设计、评估与快速迭代]
\textbf{【任务情境】}：在一个由软件架构师、前端开发人员、产品经理、UI/UX 设计师及行业临床专家组成的敏捷跨职能团队中，协同打造面向复杂工作流的移动端协作系统。\\
\textbf{【能力陈述】}：系统整合应用人机交互各项知识单元，跨越专业术语壁垒进行高效技术沟通；将模糊的用户期望转化为明确的技术约束与接口规范，依托敏捷反馈闭环推动软件增量开发与持续重构，切实解决现实工业级难题。\\
\textbf{【能力领域】}：软件与应用工程 \quad \textbf{【能力单元】}：需求工程、系统设计、系统开发、技术文档、系统评估、项目管理、用户接受度、系统改进\\
\textbf{【核心要求】}：计算机核心 \quad \textbf{【达标等级】}：\textbf{开发}
\end{taskbox}

\begin{taskbox}[任务 5：基于高层设计原型编制高精度用户界面工程开发规格说明书]
\textbf{【任务情境】}：交互设计团队交付了业务系统高保真原型与设计规范系统，工程研发团队需要将其转化为严格、无歧义的前端与移动端软件工程实施规范。\\
\textbf{【能力陈述】}：紧扣可用性标准、用户真实需求与既定设计方案，编写详尽的软件用户界面需求规格说明书（SRS），精准界定页面状态机迁移、组件封装参数、网络弱网降级逻辑、响应式布局断点及前后端接口通信契约。\\
\textbf{【能力领域】}：软件与应用工程 \quad \textbf{【能力单元】}：需求工程、技术文档编制\\
\textbf{【核心要求】}：计算机核心 \quad \textbf{【达标等级】}：\textbf{开发}
\end{taskbox}

\begin{taskbox}[任务 6：支撑全球化出海业务系统的多语言与跨文化国际化 (i18n) 架构重构]
\textbf{【任务情境】}：某企业原本面向本土市场的企业管理与协同办公软件拟拓展至中东、欧洲等多国市场，系统出现布局溢出崩溃、文字截断、阅读流向错误与文化禁忌符号等严重问题。\\
\textbf{【能力陈述】}：系统化评估现有软件代码与界面实现，准确定位受文化、语言和地域习俗影响的硬编码文本、图标隐喻、日期/货币格式以及从右至左（RTL）排版依赖，提出全面的国际化与本地化架构重构规划。\\
\textbf{【能力领域】}：软件/系统/应用工程 \quad \textbf{【能力单元】}：系统设计、系统评估、社会议题调适、系统持续改进\\
\textbf{【核心要求】}：计算机核心 \quad \textbf{【达标等级】}：\textbf{评估}
\end{taskbox}

\begin{taskbox}[任务 7：实施形成性用户调研与实验，为全新交互系统提供战略级设计演进建议]
\textbf{【任务情境】}：某科技公司拟立项研发一款基于生成式大模型的专业程序员辅助编程助手，处于概念预研期，需摸清开发者的编码心智模型与核心交互痛点。\\
\textbf{【能力陈述】}：科学设计、组织并执行形成性用户体验研究（情境调查、出声思考走查、深度访谈），精准辨识交互原型的核心优势、认知阻碍与工程痛点，提炼出驱动系统架构长效演进的高价值设计策略建议。\\
\textbf{【能力领域】}：应用工程与理论计算 \quad \textbf{【能力单元】}：方案设计、测试验证、用户接受度评测、系统改进\\
\textbf{【核心要求】}：计算机核心 \quad \textbf{【达标等级】}：\textbf{评估}
\end{taskbox}

\begin{taskbox}[任务 8：跨部门解读、提炼并传达多源用户评估反馈，驱动研发版本敏捷迭代]
\textbf{【任务情境】}：某大型 SaaS 平台在灰度内测阶段，来自企业 IT 主管、一线业务操作人员、客服工单及系统埋点日志的多源反馈出现严重分歧与术语隔阂。\\
\textbf{【能力陈述】}：科学搭建反馈收集矩阵；结合不同角色评估者的立场视角、专业语境与优先级诉求，对海量质性与量化反馈进行系统归纳与根因溯源，撰写层次分明的技术摘要，制定指导软件敏捷迭代的整改路线图。\\
\textbf{【能力领域】}：软件工程 \quad \textbf{【能力单元】}：需求工程、测试验证、技术文档、用户接受度、社会议题调适、系统改进\\
\textbf{【核心要求】}：计算机核心 \quad \textbf{【达标等级】}：\textbf{评估}
\end{taskbox}

\begin{taskbox}[任务 9：面向特定敏感脆弱人群部署交互系统的社会伦理与技术风险综合审计]
\textbf{【任务情境】}：某智慧校园项目拟在中小学课堂全面部署基于人脸微表情识别与注视点追踪的注意力实时监测与情绪预警系统。\\
\textbf{【能力陈述】}：深入调研未成年人等特殊脆弱受众群体的身心发展特征；从数据隐私泄露、知情同意权缺位、心理胁迫焦虑、算法偏见与权力结构不对等维度，全面审计该系统潜在的多维负面社会冲击，出具严密的伦理风险审查与工程叫停/整改报告。\\
\textbf{【能力领域】}：应用工程 \quad \textbf{【能力单元】}：需求工程、系统评估、用户接受度审查、社会议题调适、系统改进\\
\textbf{【核心要求】}：计算机核心 \quad \textbf{【达标等级】}：\textbf{评估}
\end{taskbox}

\begin{taskbox}[任务 10：遵循工业级规范标准，端到端高保真构建可交互用户界面原型]
\textbf{【任务情境】}：根据产品规格说明书与主流工业级设计规范（如 Material Design / Apple HIG），为某物联网智能中控面板独立开发可交互、高保真的全功能软件界面原型。\\
\textbf{【能力陈述】}：严格依据软件设计规格书与专业工程师行为规范，熟练运用现代组件库、响应式布局引擎与状态管理模型，开发出视觉一致、交互灵敏、支持无障碍遍历且达到生产级工业标准的高保真交互原型系统。\\
\textbf{【能力领域】}：软件、应用与理论工程 \quad \textbf{【能力单元】}：需求转化、交互设计、系统开发\\
\textbf{【核心要求】}：计算机核心 \quad \textbf{【达标等级】}：\textbf{开发}
\end{taskbox}

\section{课程设计与开课规划建议}

联合工作组面向各高校计算机科学专业主修、辅修及跨学科方向，推荐三套典型示范课程方案：

\subsection{计算机专业主修与辅修核心课：《人机交互导论》（建议 30 学时）}
\begin{itemize}
    \item \textbf{学时结构分配}：
    \begin{itemize}
        \item 人机交互/理解用户 （单元 1）：7 学时
        \item 人机交互/设计中的问责与责任 （单元 2）：2 学时
        \item 人机交互/无障碍与包容性设计 （单元 3）：4 学时
        \item 人机交互/交互设计评估 （单元 4）：3 学时
        \item 人机交互/交互系统设计 （单元 5）：10 学时
        \item 人机交互/社会伦理交叉贯穿 （单元 6）：2 学时
    \end{itemize}
    \item \textbf{先修要求}：敏捷软件开发基础（或基础软件工程方法）。
    \item \textbf{能力目标陈述}：修完本课程的学生应能够透彻阐述以用户为中心的设计（UCD）原则，并将其系统应用于小型工程项目的全流程设计、实现与评估中。
    \item \textbf{课程特色与教学模式}：本课程采取“项目驱动、贯穿迭代”的教学范式。前 4 周奠定扎实的交互设计基础；在第 5--6 周，学生通过一个由教师给定的短期微型项目进行热身实操，激发探索深层技术概念的学习热情；在此基础上，学生团队自主对接真实社区或特定用户群体，运用快速原型迭代法完成完整的需求分析、交互设计、工程编码与可用性评估闭环。
    \item \textbf{15 周学期教学进度大纲}：
    \begin{enumerate}
        \item 交互设计导论与核心原则 （单元 1, 5）
        \item 认知心理学视角：思考、行动与评估三元模型 （单元 1, 4, 5）
        \item 人类记忆模型与操作失误防御设计 （单元 1, 3, 4）
        \item 交互设计原则与敏捷过程模型 （单元 5）
        \item[5--6.] 设计流程与现代软件工程开发的深度融合实战 （单元 5）
        \item[7.] 设计思维与启发式评估实战 （单元 1, 4）
        \item[8.] 无障碍设计规范（WCAG）与包容性技术 （单元 2, 3, 6）
        \item[9.] 视觉界面设计美学与用户画像进阶 （单元 1, 3, 4, 5）
        \item[10.] 期末大项目：用户同理心调研与需求精准界定 （单元 1～6）
        \item[11.] 期末大项目：头脑风暴构思与低保真原型制作 （单元 1, 4, 5）
        \item[12--13.] 期末大项目：交互系统工程实现与编码 （单元 5）
        \item[14.] 期末大项目：受试者可用性实测与量化评测 （单元 4）
        \item[15.] 期末大项目：成果路演汇报与评估文档编制 （单元 1～6）
    \end{enumerate}
\end{itemize}

\subsection{跨学科特色核心课：《数据可视化导论》（建议 47 学时）}
\begin{itemize}
    \item \textbf{学时结构分配}：
    \begin{itemize}
        \item 图形交互/数据可视化单元：30 学时
        \item 图形交互/基础渲染技术单元：10 学时
        \item 人机交互/理解用户单元：3 学时
        \item 社会伦理/多元性、隐私与数据伦理单元：4 学时
    \end{itemize}
    \item \textbf{先修要求}：面向对象程序设计基础 (CS2)、线性代数。
    \item \textbf{能力目标陈述}：修完本课程的学生应能够科学筛选与清理复杂数据集，验证数据的真实性、代表性与适用边界；设计、开发并测试一套直观揭示数据内在规律、高度易用且交互流畅的数据可视化分析系统。
\end{itemize}

\subsection{进阶实战专业课：《可用性工程与系统评测》}
\begin{itemize}
    \item \textbf{先修要求}：人机交互导论、社会科学研究方法、应用数理统计学。
    \item \textbf{课程特色与教学模式}：本课程为完全基于真实工程项目的进阶实战课，深度聚焦工业级交互产品的正式形式化评估。课程全面涵盖：评测目标科学确立、目标受试者分层抽样招募、测试任务与情境剧本构建、实验室环境搭建、测试方案编写与执行、实验日志与行为数据分析、评估白皮书撰写与整改建议答辩。使学生掌握针对复杂软硬件系统在多元真实场景下开展可用性工程评测的专业技术、规程与标准协议。
    \item \textbf{14 个核心专题模块}：
    \begin{enumerate}
        \item 可用性评测研究方案的系统规划
        \item 核心评测目标与科学假设确立
        \item 被试群体特征界定与代表性抽样策略
        \item 典型工程测试任务选取与情境剧本设计
        \item 可用性量化度量指标体系构建
        \item 测试指导手册、保密协议与知情同意书编制
        \item 测试软硬件设备配置与多机位实验室搭建
        \item 预实验与方案校准
        \item 正式可用性测试的主持与出声思考法引导
        \item 行为日志与量表数据制表统计与质性扎根编码
        \item 提出优先级分明的界面重构与代码整改建议
        \item 评估结论的专业化可视化沟通与呈现
        \item 制作典型操作阻塞点“高光剪辑录像”
        \item 推动产品架构持续演进与研发流程变革
    \end{enumerate}
    \item \textbf{核心能力产出}：
    \begin{itemize}
        \item 独立设计严密规范的工业级可用性测试方案；
        \item 依据严格抽样标准招募并管理多元代表性受试者；
        \item 熟练主持并执行无偏倚的高质量可用性测试会话；
        \item 科学统计实验数据，深度剖析根因并提出可辩护的重构方案；
        \item 面向工程团队与管理层清晰汇报评估发现与商业洞察；
        \item 撰写完全符合行业规范的软件可用性评估与改进工程白皮书。
    \end{itemize}
\end{itemize}

\section{委员会专家名单}

\begin{itemize}
    \item \textbf{主席}：
    \begin{itemize}
        \item Susan L. Epstein (纽约市立大学亨特学院及研究生中心，美国)
    \end{itemize}
    \item \textbf{委员}：
    \begin{itemize}
        \item Sherif Aly (开罗美国大学，埃及)
        \item Jeremiah Blanchard (佛罗里达大学，美国)
        \item Zoya Bylinskii (Adobe 研发实验室，美国)
        \item Paul Gestwicki (鲍尔州立大学，美国)
        \item Susan Reiser (北卡罗来纳大学阿什维尔分校，美国)
        \item Amanda M. Holland-Minkley (华盛顿与杰斐逊学院，美国)
        \item Ajit Narayanan (谷歌公司，美国)
        \item Nathalie Riche (微软全球研究院，美国)
        \item Kristen Shinohara (罗切斯特理工学院，美国)
        \item Olivier St-Cyr (多伦多大学，加拿大)
    \end{itemize}
\end{itemize}

\clearpage
\onecolumn
\section{附录：核心主题与技能要求等级对照表}

下表系统梳理人机交互知识领域各知识单元的核心知识点、认知等级（Bloom 分类）、知识范畴（计算机核心 / 领域核心）及推荐学时分配：

\begin{longtable}{p{0.18\textwidth} p{0.48\textwidth} c c c}
\toprule
\textbf{知识单元} & \textbf{核心知识点与具体主题} & \textbf{认知等级} & \textbf{分类} & \textbf{学时} \\
\midrule
\endfirsthead
\toprule
\textbf{知识单元} & \textbf{核心知识点与具体主题} & \textbf{认知等级} & \textbf{分类} & \textbf{学时} \\
\midrule
\endhead
\bottomrule
\endfoot
\multirow{4}{*}{\makecell[l]{人机交互/理解用户}} 
& 用户中心设计哲学（“你不是用户”）、需求挖掘、访谈与问卷调研 & 应用 & 计算机核心 & 2 \\
& 用户画像图谱、敏捷用户故事、同理心地图与用户体验旅程图构建 & 应用 & 领域核心 & 2 \\
& 用户的生理感知、认知注意力机制、记忆模型与操作认知偏差 & 解释 & 领域核心 & 2 \\
& 多元用户群体适配、国际化与跨文化设计、多用户群体协同通信计算 & 评估 & 领域核心 & 1 \\
\midrule
\multirow{2}{*}{\makecell[l]{人机交互/设计中的\\问责与责任}} 
& 系统设计影响（绿色可持续性、社会包容性、安全、隐私与歧视防范） & 评估 & 计算机核心 & 2 \\
& 价值敏感设计（VSD）、多方利益相关者体系、正向引导与防暗黑模式 & 评估 & 领域核心 & 2 \\
\midrule
\multirow{3}{*}{\makecell[l]{人机交互/无障碍与\\包容性设计}} 
& 无障碍法律法规、W3C WCAG 2.1/2.2 规范准则与系统级辅助技术支持 & 应用 & 计算机核心 & 2 \\
& 永久/暂时/情境障碍分类、屏幕阅读器适配与无障碍前端工程落地 & 应用 & 领域核心 & 1 \\
& 能力导向设计理念、包容性参与式研发流程与全人群数字鸿沟弥合 & 评估 & 领域核心 & 1 \\
\midrule
\multirow{2}{*}{\makecell[l]{人机交互/交互设计评估}} 
& 形成性评估与总结性评估辨析、可用性四大核心指标与伦理审查规范 & 解释 & 计算机核心 & 1 \\
& 质性编码主题分析、出声思考走查、启发式评估、对比实验与统计检验 & 评估 & 领域核心 & 2 \\
\midrule
\multirow{3}{*}{\makecell[l]{人机交互/交互系统设计}} 
& 低保真与快速原型制作技术、敏捷迭代设计闭环与多端硬件约束权衡 & 应用 & 计算机核心 & 1 \\
& 交互架构模式（MVC/MVVM/声明式 UI）、双钻设计流程、格式塔视觉法则 & 应用 & 领域核心 & 3 \\
& 多模态交互输入技术（手势/语音/触觉）、交互容错机制与优雅降级 & 开发 & 领域核心 & 2 \\
\midrule
\multirow{1}{*}{\makecell[l]{人机交互/社会伦理交叉}} 
& 通用设计人本伦理、全流程参与式包容研发与长效社会外部性审计 & 评估 & 领域核心 & 贯穿渗透 \\
\end{longtable}
